\section{Experteninterview mit Prof.~Dr.~Peter Riegler}
\label{sec:AA-01_riegler}

Dieser Abschnitt dokumentiert das leitfadengestützte Experteninterview mit Prof.~Dr.~Peter Riegler, das ergänzend zur empirischen Hauptstudie durchgeführt wurde. Ziel war es, externe fachliche Perspektiven zur Nutzung, Einführung und Akzeptanz von Wikis in organisationalen Kontexten zu erfassen und mit den Ergebnissen der eigenen Untersuchung zu triangulieren.

\subsection{Interviewkontext und Person}

Das Interview fand am 26.~September 2025 im Rahmen einer Microsoft-Teams-Sitzung statt und dauerte ca.~60 Minuten. Prof.~Dr.~Peter Riegler ist Professor für Informatik an der Ostfalia Hochschule sowie wissenschaftlicher Leiter im Bereich Hochschuldidaktik (BayZiel). Er ist Initiator und Betreiber eines Community-Wikis im Kontext der \glqq Decoding the Disciplines\grqq{}-Bewegung und verfügt über langjährige Erfahrung in der Einführung und Weiterentwicklung kollaborativer Wissensplattformen.

\subsection{Zentrale Aussagen}

\paragraph{Typische Nutzungshürden}
Riegler beschreibt insbesondere psychologische Barrieren bei weniger technikaffinen Zielgruppen. Nutzende stellen sich häufig Fragen wie \glqq Darf ich Inhalte verändern?\grqq{} oder empfinden Unsicherheit gegenüber der eigenen Kompetenz. Diese Hemmschwellen seien kulturell und organisational geprägt und würden durch unübersichtliche UX zusätzlich verstärkt.

\paragraph{Struktur und Navigation}
Als zentrale Voraussetzung für Akzeptanz benennt er eine leistungsfähige Suchfunktion sowie eine klare, kontextbezogene Verlinkung von Inhalten. Inhalte sollten nicht isoliert präsentiert werden, sondern in nachvollziehbaren Nutzungskontexten eingebettet sein.

\paragraph{Visuelle Gestaltung}
Eine konsistente visuelle Gestaltung mit verständlichen Icons, klaren Farbsystemen und stabilen Layoutmustern fördere Orientierung, Vertrauen und Nutzungsbereitschaft, insbesondere bei neuen oder unsicheren Nutzenden.

\paragraph{Didaktische Prinzipien}
Riegler verweist auf die Methode \glqq Decoding the Disciplines\grqq{}, mit der implizites Expertenwissen sichtbar gemacht werden kann. Inhalte sollten nicht nur Fakten abbilden, sondern Denkprozesse transparent machen und typische Lernhürden adressieren.

\paragraph{Einführung und Weiterentwicklung}
Er betont die Bedeutung einer aktiven Community, klarer Verantwortlichkeiten sowie iterativer Weiterentwicklung. Häufige Fehler bei der Einführung seien unzureichende Kommunikation des Nutzens, fehlende Pflegeverantwortung und mangelnde Feedbackkanäle.

\subsection{Einordnung für die vorliegende Arbeit}

Die Aussagen bestätigen zentrale Befunde der eigenen Untersuchung, insbesondere:
\begin{itemize}
    \item die Bedeutung klarer Struktur und Suchunterstützung,
    \item die Rolle von Kommunikation und Vorbildern für Akzeptanz,
    \item die Notwendigkeit organisatorischer Governance-Strukturen,
    \item die Relevanz impliziter UX-Qualitäten für Vertrauen und Nutzung.
\end{itemize}

Das Interview dient somit als externe Validierung der empirischen Ergebnisse.
